
\subsubsection{6.1 Key Findings and Implications}

\textbf{Trustworthiness rankings are stable latent model properties.} High partial $\rho$ for both fairness and adversarial robustness after capability control indicates that trustworthiness-dimension rankings are not test-specific artifacts. For practitioners: model selection decisions based on in-distribution trustworthiness rankings are more robust to distribution shift than one might assume from the adversarial benchmarking literature. The relative ordering of models on fairness and robustness benchmarks is preserved when those benchmarks become harder or shift context.

\textbf{The adversarial disruption hypothesis does not generalize to diverse model populations.} The DecodingTrust N = 2 observation (GPT-4 more vulnerable to adversarial jailbreaking than GPT-3.5 on standard metrics) motivated the hypothesis. At N = 13 with diverse model families, the pattern does not scale. Both adversarial pairs show substantially positive partial $\rho$ with zero rank reversals. The models most robust to adversarial perturbation under easier conditions remain the most robust under harder conditions.

\textbf{Capability control reveals hidden differential structure.} Raw $\rho$ is near-ceiling across all dimensions ($\Delta\rho_{raw} \approx -0.005)$, making dimensions appear equally and perfectly stable. The partial analysis reveals a meaningful differential ($\Delta\rho_{partial}$ = 0.192) that is directionally consistent with the hypothesis, though below the preregistered magnitude threshold. Evaluators reporting only raw cross-benchmark correlations miss this dimension-specific information.

\textbf{The fairness-robustness differential is directional but unconfirmed.} Fisher z p = 0.024 supports the directional hypothesis. The effect falls short of the preregistered criterion by 0.008. Whether this directional advantage reflects model-internal mechanisms (stable latent bias for fairness versus partially disrupted latent properties for robustness) or benchmark structural factors (BBQ split design) cannot be determined from the current data.

\subsubsection{6.2 Limitations}

\textbf{GLUE/AdvGLUE arm structurally unavailable.} GLUE-X \citep{Yang2023GLUEX} evaluates encoder-only PLMs, with zero model overlap with the decoder-only LLMs in TrustLLM. The OOD robustness dimension from TrustLLM serves as a proxy. Direct evaluation of decoder-only LLMs on AdvGLUE would require new data collection.

\textbf{$\Delta\rho$ = 0.192 falls below the preregistered threshold.} The 0.008 shortfall is within measurement uncertainty at N = 13. Contributing factors include reduced power from three missing models and an $\rho_{ANLI}$ (0.684) higher than the mechanism hypothesis predicted. The directional claim is supported; the preregistered magnitude confirmation is not.

\textbf{Causal mechanism unknown.} The adversarial disruption mechanism is falsified. Two alternatives remain viable but untested: (1) \textit{Stable latent bias hypothesis} --- fairness-relevant stereotypical associations are more deeply encoded during pretraining than robustness-relevant properties. (2) \textit{Benchmark overlap hypothesis} --- BBQ disambiguated and ambiguous splits share sufficient item structure that the high $\rho$ reflects benchmark design rather than model mechanism. Item-level Jaccard analysis would distinguish these.

\textbf{Sensitivity to covariate choice.} Under Winogrande as alternative covariate, $\Delta\rho$ collapses to 0.073 (from 0.192 under MMLU). Robustness $\rho$ values appear more sensitive to covariate choice than fairness $\rho$ values. The $\Delta\rho$ finding is MMLU-specific and should be interpreted accordingly; richer capability covariates (BIG-Bench, GSM8K) would clarify robustness of the differential.

\textbf{Scope restricted to 2023--2024 model population (N = 16).} Results pertain to the specific model families evaluated by TrustLLM. Whether post-2024 models exhibit different stability patterns is not addressed.

\textbf{P3 prediction (instruction-tuning effects) untested.} Whether RLHF-aligned models show greater differential fairness generalization compared to base models was not examined. This remains the most practically actionable open question.

\textbf{N asymmetry across dimensions.} Fairness analysis uses N = 16; robustness uses N = 13. The $\Delta\rho$ comparison is conducted on the N = 13 intersection to satisfy the Meng et al. [1992] same-sample requirement, but the fairness estimate on N = 13 ($\rho$ = 0.967) differs slightly from the full-sample estimate ($\rho$ = 0.962). The N = 16 result is the primary fairness finding; N = 13 is used only for the differential test.

\subsubsection{6.3 Broader Impact}

This study provides empirical grounding for an implicit assumption in LLM safety evaluation. The finding that trustworthiness rankings are stable supports benchmark-based model selection as a meaningful signal for deployment decisions. However, rank stability and absolute performance adequacy are distinct: models may maintain their relative ordering on harder benchmarks while all showing degraded performance. Stability of rank does not imply that any model meets an absolute safety or fairness threshold for a given deployment context.

